\section*{Chapter \ref{ch:rc:short-rate-models} --- Short-Rate Models}

\begin{solution}{exo:rc:short-rate-models:1}
$\sigma B(\tau)/\tau$: at ten years $80\times(1-e^{-0.3})/0.3 = 69.1$ basis points; at
thirty years $80\times(1-e^{-0.9})/0.9 = 52.7$.
\end{solution}

\begin{solution}{exo:rc:short-rate-models:2}
Its bond prices depend on three constants ($\kappa$, $\bar x$, $\sigma$) and today's short
rate, a three-parameter family of curves that cannot match an arbitrary curve.
Hull--White makes the drift's level a function of time, $\vartheta(t)$, chosen so
that every discount factor is repriced; in the Gaussian form it is $r_t =
f(0,t)+x_t$.
\end{solution}

\begin{solution}{exo:rc:short-rate-models:3}
$\bar x-\sigma^2/(2\kappa^2) = 3\%-0.0001/0.02 = 2.50\%$.
\end{solution}

\begin{solution}{exo:rc:short-rate-models:4}
$j_{\max} = \lceil0.184/(0.03/12)\rceil = 74$: 149 nodes. With weekly steps
$\lceil0.184/(0.03/52)\rceil = 319$: 639 nodes. Small $\kappa$ makes wide trees; the
width is reached only after $j_{\max}$ steps, so short trees stay narrower.
\end{solution}

\begin{solution}{exo:rc:short-rate-models:5}
It needs every bond price at expiry to be a monotone function of one state
variable, so that one critical value $x^\star$ splits the exercise region. With
two factors the exercise boundary is a curve in the $(x,y)$ plane, and the
coupon bond's exercise region does not decompose into zero-coupon regions with
a common boundary.
\end{solution}

\begin{solution}{exo:rc:short-rate-models:6}
Steeper: stronger mean reversion damps the volatility of long rates more, so
to reproduce later expiries into shorter swaps (whose volatility the market
keeps high) $\sigma(t)$ must rise faster; with a larger $\kappa$ the fit can even require
large late volatilities. $\kappa$ and $\sigma(t)$ trade off: the term structure of
volatility is shared between them.
\end{solution}

\begin{solution}{exo:rc:short-rate-models:7}
Errors against 0.026018: $+1.0\times10^{-4}$ (monthly), $+6.1\times10^{-5}$
(fortnightly), $+2.6\times10^{-5}$ (weekly), $-0.2\times10^{-5}$ (half-weekly). Halving
the step roughly halves the error until it changes sign: first order, as
expected for a tree whose payoff has a kink at the strike.
\end{solution}

\begin{solution}{exo:rc:short-rate-models:8}
Repricing co-terminals fixes the volatility of each swap rate; a \omterm{def:rc:convexity-adjustments-and-constant-maturity-products:spread}{CMS spread
option} also needs the correlation between the ten-year and the two-year rate,
which is one in any one-factor model: the spread's volatility collapses and
the option is badly underpriced. Use a two-factor model, or a copula of the
two rates' smiles, calibrated to a correlation.
\end{solution}

\begin{solution}{pb:rc:short-rate-models:1}
\textbf{1.} The co-terminals into year ten: they are the Europeans the Bermudan can
be exercised into, and its value lies between the largest of them and a price
that depends on their joint dynamics.
\textbf{2.} 60.8 to 91.0 basis points, rising with expiry.
\textbf{3.} That the rate volatility expected later (for shorter remaining swaps) is
higher than mean reversion with a flat $\sigma$ would give.
\textbf{4.} The model reprices it by construction through $\vartheta(t)$.
\textbf{5.} It is poorly identified from one column of co-terminals; it trades off
with $\sigma(t)$. Desks fix it by convention or from the relation between
co-terminal and other swaption volatilities, and keep it stable.
\textbf{6.} 86.0 basis points.
\textbf{7.} 76.0 to 76.7 basis points.
\textbf{8.} $+15.3$ basis points at one year, $-14.2$ at nine years.
\textbf{9.} Early exercises are overpriced (the model's short expiries are too
volatile), late ones underpriced.
\textbf{10.} With $\kappa=3\%$ and constant $\sigma$, co-terminal volatilities are nearly flat in
expiry, while the market's rise by 30 basis points.
\textbf{11.} 68.7 basis points in the first year, 116.3 at the highest (years seven
to eight).
\textbf{12.} The last co-terminal (nine into one) has a volatility of 91.0 only
slightly above the eight-year's; to reprice it with the earlier levels held,
the last interval needs less.
\textbf{13.} Exactly the co-terminals; other swaptions only as far as $\kappa$ and the
fitted $\sigma(t)$ happen to reproduce them.
\textbf{14.} That it stays at the last level, 104.2 basis points.
\textbf{15.} Noisy parameters and P\&L jumps from day to day; hedge ratios that change
with the calibration rather than the market. Smooth or regularise, and monitor
the parameters.
\textbf{16.} Hull--White with $\kappa$ fixed at 3\% and a piecewise-constant $\sigma(t)$ on the
co-terminal expiries, fitted daily to the co-terminal at-the-money swaptions;
its limits: one factor, flat smile, no fit to other swaptions.
\textbf{17.} Show the co-terminal fit errors with one and with nine volatilities, the
stability of $\sigma(t)$ over past months, and the Bermudan's price under both,
against a benchmark model (a market model, chapter~8).
\textbf{18.} Imperfect correlation between the rates exercised into, which lowers the
switch value between exercise dates; the cost is slower pricing, more parameters
and a correlation to mark.
\textbf{19.} Named result: \emph{one sigma is not enough}: the best constant volatility
misses the co-terminals by $+15.3$ basis points at one year and $-14.2$ at nine;
the fitted $\sigma(t)$ runs from 68.7 to 116.3 basis points.
\textbf{20.} The curve exactly and the volatilities of the options the product can
become.
\end{solution}

\begin{solution}{iq:rc:short-rate-models:1}
It pulls the short rate back towards its path, so shocks to it die out and
long rates move less than short ones: it sets how volatility decreases with
maturity ($\sigma B(\tau)/\tau$), and how correlated the exercises of a Bermudan are.
A trader reads it as the model's slope of volatility across tenors.
\iqlookfor{mean reversion as a volatility-shape parameter.}
\end{solution}

\begin{solution}{iq:rc:short-rate-models:2}
$r$ is Ornstein--Uhlenbeck; $\int_t^Tr_s\,ds$ given $r_t$ is normal with mean $\bar
x\tau+(r_t-\bar x)B(\tau)$ and variance $\frac{\sigma^2}{\kappa^2}(\tau-B-\frac\kappa2B^2)$; so
$P = \exp(-\text{mean}+\text{variance}/2)$, an exponential-affine function of $r_t$.
\iqlookfor{Gaussian integral and the affine form.}
\end{solution}

\begin{solution}{iq:rc:short-rate-models:3}
Build the tree for the zero-mean state first; then step forward, carrying the
Arrow--Debreu prices of each node; at each step choose the shift $\alpha_i$ so that
the sum of Arrow--Debreu prices times one-step discount factors equals the
curve's discount factor to the next date; update the Arrow--Debreu prices with
the branching probabilities and the shifted rates.
\iqlookfor{forward induction with Arrow--Debreu prices.}
\end{solution}

\begin{solution}{iq:rc:short-rate-models:4}
The co-terminal swaptions (each exercise date into the swap to the final
date), because the Bermudan is exactly one of them once exercised; a one-factor
model cannot fit the whole matrix, and fitting instruments the product never
becomes spends parameters on irrelevant risks.
\iqlookfor{the replication logic of the calibration set.}
\end{solution}

\begin{solution}{iq:rc:short-rate-models:5}
Products depending on imperfect correlation between rates: spread options,
curve-steepener exotics, and to a lesser degree Bermudans (switch value). Test
by pricing with a two-factor model calibrated to the same instruments and
measuring the difference; if it matters, use the richer model or reserve.
\iqlookfor{correlation, and a benchmark to measure it.}
\end{solution}

\begin{solution}{iq:rc:short-rate-models:6}
Check the tree reprices discount factors (the $\alpha$ fit); halve the step and see
if the error halves (discretisation) or stays (a bug); compare zero-coupon bond
options with the closed form; check the exercise date and schedule alignment
with the tree's grid, the branching at the edges and the probabilities (all in
$[0,1]$, summing to one).
\iqlookfor{convergence study and component tests.}
\end{solution}
